\chapter{Web-Based Evaluation and Validation Loop}\label{sec:web}

Chapter~\ref{sec:eval} showed that the shared representation is trainable on both construction paths. This chapter closes the validation chain by showing how a trained, representation-compatible ETA predictor is used inside the \WebApp{} to produce persistent prediction--outcome evidence. The web application is not presented as a production navigation service; it is an evaluation harness that connects route selection, model inference, simulated or replayed movement, realized arrival, and aggregate analysis.

\section{Role of the web loop}
The web loop validates operational usability. A trained model can have good held-out metrics, yet still be difficult to use in an interactive setting if the surrounding system cannot load the network, compute a route, invoke inference, track the journey, or preserve the result. The \WebApp{} exercises these steps in one workflow:
\begin{enumerate}
    \item load a representation-compatible bundle and render the road network;
    \item let the user select origin and destination points;
    \item compute a SUMO shortest route for the selected endpoints;
    \item invoke the attached ETA model and display the predicted arrival;
    \item run the journey until the vehicle reaches its destination;
    \item compare predicted and realized arrival times;
    \item store the completed journey for later analysis.
\end{enumerate}

The canonical screenshots for this process appear in Section~\ref{sec:system:web:ui}. Figure~\ref{fig:web_ui_init} shows the loaded dashboard, Figure~\ref{fig:web_ui_route} shows route selection and journey start, and Figure~\ref{fig:web_ui_recent} shows the completed prediction--outcome record.

\section{Per-journey evidence}
The atomic evidence unit is one completed journey. Each journey record contains the selected route, trip distance, start time, predicted ETA, realized arrival, actual duration, signed error, absolute error, and accuracy. This is the bridge between model training and user-facing validation: the model produces an ETA before the trip is complete, and the application later records what actually happened.

Journey \#2709 in Figure~\ref{fig:web_ui_route}--\ref{fig:web_ui_recent} is the running example. The selected route is 13.51\,km. The trip starts on Wednesday at 05:21:00, with a predicted arrival at 05:38:11, corresponding to a predicted duration of 17:11 minutes. At completion, the recorded prediction is 19 seconds before the actual arrival, giving 98.1\% journey-level accuracy. This example is not used as a statistical claim by itself; it illustrates the evidence schema that can be accumulated over many journeys.

\section{Path-independent validation}
The same web surface is used regardless of how the underlying graph bundle was produced. The simulation-backed workflow uses controlled traffic from the desktop simulation path, while Figure~\ref{fig:web_ui_porto} shows the same client operating on the Porto trajectory-conversion network. This matters because it confirms that path independence extends beyond file layout and model training: once exported, both bundles can support the same route-selection and ETA-evaluation workflow.

\section{Aggregation and reporting}
Completed journeys are persisted so that the interface can compute aggregate error panels, recent-journey summaries, plots, and report exports without rerunning inference. Aggregates such as MAE, RMSE, and MAPE are meaningful here because each row has both a prediction and a realized outcome. The dashboard therefore serves two roles: it lets an operator inspect single journeys, and it accumulates enough evidence to analyze cohorts by route, duration, time of day, or dataset source.

\section{Validation boundary}
The web loop validates that the representation, trained model, backend service, and user-facing client can work together end to end. It does not claim universal production readiness, optimal routing, or model superiority across all possible cities and traffic regimes. Those claims would require broader deployments and additional controlled comparisons. The contribution here is narrower and directly tied to the report: datasets produced by both construction paths can train a compatible ETA model, and that trained model can be exercised in a persistent interactive validation loop.

\FloatBarrier
